\documentclass{article}
\usepackage{amsmath}
\usepackage{amssymb}
\usepackage{graphicx}
\usepackage[utf8]{inputenc}

\title{Sound Localising Methods}
\date{}

\begin{document}

\maketitle

\section*{1. Model the incoming sound}

First we have to model the incoming sound somehow.

\begin{equation}
    x_i(t) = s(t) * h_i(t) + n_i(t)
\end{equation}

Where:
\begin{itemize}
    \item $x_i(t)$ is the signal captured by the $i$-th microphone.
    \item $s(t)$ is the source signal.
    \item $h_i(t)$ is the room response. It consists of the time difference that the sound has to travel and all other reverberations. (Convolution can apply these features at the same time).
    \item $n_i(t)$ is the noise captured by the $i$-th microphone.
\end{itemize}

\section*{2. Cross-correlation}

Then basically we'll calculate cross-correlation (later we'll use tricks but the core is still a cross-correlation problem).

So:
\begin{equation}
    R_{ij}(\tau) = \int_{0}^{T} x_i(t) x_j(t+\tau) \, dt
\end{equation}

In theory we could do this but when we're working with computers we should  apply discrete methods.

\begin{equation}
    R_{ij}[\tau] = \sum_{k=0}^{N} x_i[k] x_j[k+\tau]
\end{equation}

It is a cross correlation that gives us back how the two signals match as a function of $\tau$, then we can select the $\tau$ that maximizes this function and say that's the time difference of arrival.

\section*{3. Frequency Domain}

But we can do it in a way more sophisticated way by looking at these events in frequency domain.

But before we do that it's worth noticing that the computation of cross-correlation and convolution is really similar. And convolution in time domain becomes multiplication in frequency domain.

So:
\begin{equation}
    R_{ij} = x_i(t) \star x_j(-t)
\end{equation}
(The minus sign becomes complex conjugation).

\begin{equation}
    G_{ij}(\omega) = X_i(\omega) X_j^*(\omega)
\end{equation}

Before we transform back to get the TDOA we whiten the signals by removing the magnitude. This makes the method more precise.

So we'll transform $\frac{G_{ij}}{|G_{ij}|}$ back to time domain.

\begin{equation}
    R_{ij}^{PHAT}(\tau) = \frac{1}{2\pi} \int_{-\infty}^{\infty} \frac{X_i(\omega) X_j^*(\omega)}{|X_i(\omega) X_j^*(\omega)|} e^{i\omega\tau} \, d\omega
\end{equation}

So by doing this we'll get a pulse at the estimated TDOA.

\vspace{1em}
\noindent \textbf{So with GCC-PHAT we can get the time delays for every pair of mics.}

\end{document}
